\documentclass[book.tex]{subfiles}

\begin{document}


\chapter{Special Relativity}

\label{chap:special_relativity}
\noindent\rule{11cm}{0.4pt} \\
\textbf{Prerequisites:} \autoref{chap:lorentz_group} \autoref{chap:newtons_laws} \\
\textbf{Difficulty Level:} ** \\
\noindent\rule{11cm}{0.4pt}
\vspace{5mm}

Previously in physical theories we consiered, we limited the range of speeds addressable by a physical theory, blaming relativistic effects. What precisely is a relativistic effect? Broadly, relativity addresses the problem of changing reference frames correctly in physics. 

Here, a reference frame is a point in spacetime (typically modeled as $\mathbb{R}^4$) and denoted by 
\begin{align}
(t, x, y, z) \in \mathbb{R}^4
\end{align} 
By convention, the first parameter denotes the time, and the remaining parameters the spatial coordinates. Each reference frame is also attached to a basis for $\mathbb{R}^4$ called a coordinate system.

You can think of the reference frame and coordinate system as hidden parameters in all the models we've dealt with so far. Often, it is useful to switch between two reference frames. Let's consider an example of a change of reference frame in Figure \ref{fig:standard_conf}.

\begin{figure}[ht]
   % \includegraphics[width=0.7\linewidth]{figures/Physics/mechanics/special_relativity/standard_conf.png}
   \includegraphics[width=0.8\linewidth]{figures/Physics/mechanics/special_relativity/Changing ref frame.png}
   \caption{Changing the reference frame.}
\label{fig:standard_conf}
\end{figure}
% TODO: This figure is sourced from \url{https://en.wikipedia.org/wiki/File:Standard_conf.png}. Replace with our own illustration.

This figure corresponds to the transformation 
\begin{align}
    x' &= x - vt \\
    y' &= y \\
    z' &= z \\
    t' &= t
\end{align}
This type of transformation is typically called a uniform motion. In this case, the uniform motion is purely in the $x$ coordinate direction. There are in general 3 types of transformations allowed in classical physics.

%\textcolor{red}{TODO: Add section on orthogonal transformations to math chapter.}
%\textcolor{red}{TODO: How do I obtain the new coordinate basis $S'$ for each of these transformations.}


\begin{itemize}
    \item Uniform Motion: A uniform motion with velocity $\vec{v} \in \mathbb{R}^3$ performs the the transformation.
    \begin{equation}
        (\vec{x}, t) \to (\vec{x}+t\vec{v}, t)
    \end{equation}
    \item Translation: A translation is given by position shift $\vec{a} \in \mathbb{R}^3$ and time shift $s \in \mathbb{R}$ and performs the transformation
    \begin{equation}
        (\vec{x}, t) \to (\vec{x} + \vec{a}, t + s)
    \end{equation}
    \item Rotation: A rotation is given by an orthogonal matrix $G: \mathbb{R}^3 \to \mathbb{R}^3$ and performs the transformation
    \begin{equation}
        (\vec{x}, t) \to (G\vec{x}, t)
    \end{equation}
\end{itemize}

\section{Galilean Transformations}
Galilean transformations are the set of transformations which can be written as the compositions of uniform motions, translations, and rotations. 

The collection of Galilean transformations forms a mathematical group. You will show that this group is a Lie group of dimension 10 in the exercises.

\section{Lorentzian Transformations}
One of the great discoveries in physics was Einstein's realization that Galilean transformations were not sufficient to describe physical phenomena at high speeds. It turned out that a more general class of transformations, called Lorentz transformations were needed for this.

The Lorentz group $\mathcal{L}$ is the set of transformations $O(3, 1)$ that leave the Minkowski metric of a 4-vector unchanged. A Lorentz boost is one such transformation

\begin{align}
    x' &= \gamma(x - vt) \\
    y' &= y \\
    z' &= z \\
    t' &= \gamma \left ( t - \frac{vx}{c^2} \right )
\end{align}
where
\begin{align}
\gamma &= \frac{1}{\sqrt{1 - \frac{v^2}{c^2}}}
\end{align}
%\textcolor{red}{TODO: Introduce Lorentz Transformations and the Lorentz group.}

%\textcolor{red}{TODO: Talk about how vector fields and tensor fields change under these various transformations.}
Note that in the case that $v \ll c$, we have that $\gamma \approx 1$, in which case the Lorentzian transformation reduces to a Galilean transformation. Similar transformations can be defined in the $y$ and $z$ directions as well. In general, it is possible to define a Lorentzian boost in an arbitrary direction $n$
\begin{align}
    t' &= \gamma \left ( t - \frac{v n \cdot r}{c^2} \right ) \\
    r' &= r + (\gamma - 1)(r \cdot n)n - \gamma t v n
\end{align}
\subsection{Exercises}
\begin{enumerate}
    \item Prove that the set of Galilean transformations forms a Lie Group.
\end{enumerate}

\end{document}